\documentclass[11pt,a4paper]{article}

\usepackage[margin=2.5cm]{geometry}
\usepackage{fontspec}
\setmainfont{Latin Modern Roman}
\setsansfont{Latin Modern Sans}
\setmonofont{Hack}[Scale=0.85]
\usepackage{amssymb}
\usepackage{amsmath}
\usepackage{graphicx}
\usepackage[dvipsnames,table]{xcolor}
\usepackage{tikz}
\usetikzlibrary{positioning,arrows.meta,fit,backgrounds,shapes.geometric,calc,decorations.pathmorphing}
\usepackage{listings}
\usepackage{enumitem}
\usepackage{hyperref}
\usepackage{fancyhdr}
\usepackage{titlesec}
\usepackage{parskip}
\usepackage{float}
\usepackage{booktabs}
\usepackage{array}
\usepackage{adjustbox}
\usepackage[most]{tcolorbox}

\definecolor{router}{HTML}{E74C3C}
\definecolor{admitted}{HTML}{2ECC71}
\definecolor{denied}{HTML}{E74C3C}
\definecolor{tls}{HTML}{2ECC71}
\definecolor{cipher}{HTML}{9B59B6}
\definecolor{session}{HTML}{3498DB}
\definecolor{ca}{HTML}{F39C12}
\definecolor{codebg}{HTML}{F8F8F8}
\definecolor{codeframe}{HTML}{DDDDDD}
\definecolor{nixblue}{HTML}{5277C3}
\definecolor{soft}{HTML}{F4F6F8}

\newtcolorbox{abstractbox}[1][]{
  enhanced, colback=soft, colframe=nixblue,
  boxrule=0.6pt, arc=2pt, left=10pt, right=10pt, top=8pt, bottom=8pt,
  title=\textbf{Abstract}, coltitle=white, colbacktitle=nixblue,
  fonttitle=\sffamily\bfseries, #1
}
\newtcolorbox{findingbox}[1][]{
  enhanced, colback=soft, colframe=router,
  boxrule=0.6pt, arc=2pt, left=10pt, right=10pt, top=8pt, bottom=8pt,
  title=\textbf{Finding}, coltitle=white, colbacktitle=router,
  fonttitle=\sffamily\bfseries, #1
}

\lstset{
  basicstyle=\ttfamily\small,
  backgroundcolor=\color{codebg},
  frame=single,
  rulecolor=\color{codeframe},
  breaklines=true,
  breakatwhitespace=true,
  tabsize=2,
  showstringspaces=false,
}

\pagestyle{fancy}
\fancyhf{}
\fancyhead[L]{\small\textit{Iterating on the launcher UI and the agent graph view}}
\fancyhead[R]{\small\textit{DVerse Platform}}
\fancyfoot[C]{\thepage}
\renewcommand{\headrulewidth}{0.4pt}

\hypersetup{
  colorlinks=true,
  linkcolor=blue!60!black,
  urlcolor=blue!60!black,
}

\title{%
  \vspace{-1cm}
  {\LARGE\bfseries Iterating on the Launcher UI and the Agent Graph View:\\
  From egui Admission Window to a Stable Force-Directed Graph}
}
\author{Yordan Mitev}
\date{\today}

\begin{document}
\maketitle
\thispagestyle{fancy}

\begin{abstractbox}
  Across the semester the operator-facing surface of the DVerse
  platform moved through three shapes: a native egui admission
  window bundled into the Zenoh router binary, a Tauri desktop
  launcher that absorbed the router as a library and re-implemented
  the admission flow in React, and a live agent/topic graph that
  went through three layout attempts before settling on a
  force-directed placement with a per-topology stability cache.
  Each shape started from a hypothesis about where operator
  attention should live, and each of the first two was closed by
  something concrete that the next shape set out to fix. The graph
  iteration is the strongest disconfirmation case: the layout
  engine was not the problem, the snapshot-driven re-layout was.
\end{abstractbox}

\vspace{1em}

\begin{table}[H]
  \centering
  \small
  \begin{tabular}{|l|l|p{8cm}|}
    \hline
    \rowcolor{gray!15}
    \textbf{Version} & \textbf{Date} & \textbf{Changes} \\
    \hline
    V1 & 21 June 2026 & Initial writeup of the V1 \(\rightarrow\)
    V2 \(\rightarrow\) V3 iteration, the central
    ``snapshot-driven re-layout was the problem'' finding inside
    V3, the method and results, and the remaining open questions. \\
    \hline
  \end{tabular}
  \caption{Revision history.}
  \label{tab:ui-iter-history}
\end{table}

\newpage
\tableofcontents
\newpage

% ==============================================================
\section{Framing}
% ==============================================================

DVerse needs an operator-facing surface for two roles that overlap
on the same person at different times: the admin who decides who
joins a session, and the joiner who picks a session to join. By
sprint 2 the platform had a working mTLS fabric and an admission
protocol, but no place to drive admission from. By sprint 5 it
had a single Tauri desktop launcher hosting login, session
chooser, requesting-join wait state, the main bots view, the
admin's pending-requests panel, and a live agent/topic graph.

What follows is the trace of that move. Each iteration is a
hypothesis about where operator attention should live, an attempt
to ship it, and the finding that closed it off. The graph view
gets the largest share of space because it is the iteration
whose disconfirmations were most informative.

% ==============================================================
\section{V1: egui Admission GUI Bundled Into the Router}
% ==============================================================

\textbf{Hypothesis.} The admission decision is a property of the
router, so it belongs in a window owned by the router process. A
native eframe/egui window inside the \texttt{zenoh\_router} binary
is enough to drive admit/deny.

\textbf{What shipped.} Commit \texttt{2782235} introduced the
\texttt{zenoh\_router} crate together with the egui GUI at
\path{src/zenoh_router/src/gui.rs}. The window had a top status
bar (Starting, Running, Reloading ACL, Error), a two-column
central panel (pending CNs on the left with Admit and Deny
buttons, admitted and denied lists on the right), and a bottom
log panel that tailed the router's structured-log stream. It
polled \texttt{Arc<Mutex<AppState>{}>} every 200\,ms and queued
admit and deny actions through an \texttt{Action} enum drained by
the router loop. The file started at 109 lines and grew to
several hundred as the rest of the router's state was surfaced
through it.

\textbf{Disconfirmation.} egui locked the admission UX to whatever
machine was running the router process. Operators who were not
the admin had nothing to look at on their own machine. The
joiner side of the admission protocol (the wait for a decision
on \texttt{dverse/session/admission/<cn>}) had no rendering
anywhere. The rest of the platform's operator workflow (login,
session selection, bot management) had no obvious home that
matched the egui surface in style or in lifecycle, so each piece
risked landing somewhere different. egui was correct for the
admin role on the admin machine, and the right shape for nothing
else.

% ==============================================================
\section{V2: Tauri Launcher, Router as a Library}
% ==============================================================

This iteration is two commits sequenced through one decision: the
operator surface should be a proper cross-platform desktop app,
not an in-process egui window, and the router should be embedded
in it rather than run alongside it.

\textbf{Hypothesis.} Move the operator surface into a Tauri 2
desktop app with a React, Vite, and Tailwind frontend. Then,
once that app exists, embed the router into it as a library so
admins and joiners both run one process and look at one window.
The admission flow gets re-implemented inside the React UI where
the rest of the operator workflow lives.

\textbf{What shipped.}

\begin{itemize}[leftmargin=1.4em]
  \item \textbf{Initial Tauri launcher} (\texttt{df4cbde}, by Denis
        Lorenzo Neagoe). New tree at
        \path{experiments/zenoh-client-launcher/} with a React +
        Vite + Tailwind frontend and a Rust Tauri backend. The
        backend exposed \texttt{start\_bot}, \texttt{stop\_bot},
        and \texttt{get\_bot\_statuses} commands managing bot
        agent child processes. Initial components were
        \texttt{ConnectTab}, \texttt{BotsTab}, \texttt{BotCard},
        and \texttt{BotForm}. The launcher connected to a Zenoh
        router as a separate process. Two follow-ups (\texttt{ef3c023}
        and \texttt{44ad31b}) moved autodiscovery to the login
        screen and fixed the wiring so bots actually joined the
        network.
  \item \textbf{Router-as-library embed} (\texttt{1818090}, closes
        \#105). The \texttt{zenoh\_router} crate was split so its
        runtime lives in \path{src/zenoh_router/src/lib.rs} and
        the binary is a thin shell over it. The egui GUI was
        gated behind an \texttt{egui-gui} default feature so the
        library can build without pulling \texttt{eframe} or
        \texttt{egui}. The Tauri launcher (now at
        \path{src/zenoh_client_launcher/}) depends on
        \texttt{zenoh\_router} with
        \texttt{default-features = false} and starts
        \texttt{zenoh\_router::router::run} once at boot over a
        shared \texttt{AppState}. One binary, one window. The
        egui GUI was retired as the admin surface.
  \item \textbf{Chooser, requesting-join, pending-requests panel}
        (\texttt{45cf40f}, \#110). \texttt{ChooserScreen.tsx}
        replaced the old create-or-join radio with a credentials
        form plus a list of LAN sessions discovered through the
        existing \texttt{discover\_routers} mDNS browse, with a
        ``Create new session'' button alongside.
        \texttt{RequestingJoinScreen.tsx} renders the wait state
        between sending a JoinRequest and receiving an admission
        decision. \texttt{PendingRequestsPanel.tsx} is the
        admin's Allow / Deny list, calling
        \texttt{admit\_request} and \texttt{deny\_request}
        through the Tauri command bridge. The pending-requests
        UI now lives inside the React app, next to the bots view
        in \texttt{MainScreen}, rather than in a separate window.
\end{itemize}

\textbf{Outcome.} The operator surface collapsed into one process
and one window. Admin and joiner both run the launcher; admin
sees the pending-requests panel on the right of
\texttt{MainScreen}, joiner sees \texttt{RequestingJoinScreen}
while \texttt{join\_flow=Pending}. The egui pending-requests
pattern lived on as a UX shape (a list of pending CNs with
Allow and Deny), reimplemented in React on top of the embedded
router. The egui binary still builds (the
\texttt{egui-gui} feature is opt-in) but it is no longer the
operator surface.

% ==============================================================
\section{V3: Live Agent Graph}
% ==============================================================

V2 settled the question of where operator attention lives, but
not what the operator looks at once they get there. Beyond
``which bots are online'', the operator has no view of who
publishes what to whom. This is the gap V3 set out to close, and
it is the iteration whose disconfirmations were most informative.

\subsection{V3a: Dagre, Hierarchical Left-to-Right}

\textbf{Hypothesis.} A hierarchical left-to-right layered layout
is the right shape for ``who publishes what to whom''. Borrow
the convention from ROS2's \texttt{rqt\_graph}: agents are oval
nodes, topics are rectangles, edges are publishes and subscribes.
Use \texttt{@xyflow/react} as the renderer and \texttt{dagre} as
the layout engine.

\textbf{What shipped.} Commit \texttt{1a33dbc} introduced
\path{src/zenoh_client_launcher/frontend/src/components/GraphTab.tsx}.
The tab consumed the existing
\texttt{connected\_nodes[].agents[].publishes/subscribes} fields
of the snapshot that \texttt{App.tsx} already polled every
500\,ms. It built React Flow nodes and edges in a \texttt{useMemo}
and laid them out with \texttt{dagre} (\texttt{LR} direction).
Pure frontend, no backend change.

\textbf{Disconfirmation.} The hierarchical layout was the wrong
metaphor for a pub/sub mesh. Bidirectional relationships (a pair
of agents that publish to each other's topics) flattened into a
tangle of back-edges that defeated the layered reading. Worse,
\texttt{dagre} re-laid out the entire graph on every snapshot
tick. Status changes and heartbeat ticks (which do not change
the graph's topology at all) caused the same nodes to land in
slightly different positions, so the operator's eye lost track
of which oval was which agent between ticks.

\subsection{V3b: ELK, Layered and Orthogonal}

\textbf{Hypothesis.} A more principled layout engine will hide
the churn and route edges more cleanly. ELK's layered algorithm
with orthogonal edge routing is a known-good answer to the same
problem dagre was solving and ships better edge-bend behaviour.

\textbf{What shipped.} \texttt{cd03590} added \texttt{elkjs} as a
dependency; \texttt{4eb479c} swapped the dagre call in
\texttt{GraphTab.tsx} for an ELK layered/orthogonal layout;
\texttt{a736677} dropped \texttt{dagre} as unused.

\textbf{Disconfirmation.} ELK routed edges noticeably better,
but the underlying problem was unchanged. Nodes still jumped on
every snapshot tick, because the layout was still being recomputed
from scratch on every tick. ELK is also asynchronous, returning a
\texttt{Promise}, so a slow tick's result could land after a
fresh tick's snapshot was already rendered, producing flicker
that dagre's synchronous call had not had. The algorithm change
made one symptom better and another symptom worse, without
moving the cause.

\begin{findingbox}
The layout engine was not the problem. The pattern of
``recompute the whole layout from this tick's snapshot, then
render'' was the problem. As long as positions are a function of
the latest snapshot alone, any change to the snapshot relocates
nodes, and the human's spatial memory of which oval is which
agent breaks every 500\,ms. The fix is to make positions a
function of \emph{topology} (which agents and topics exist, and
which edges connect them) rather than of \emph{snapshot}
(everything the polling loop returns including statuses and
timestamps).
\end{findingbox}

\subsection{V3c: Force-Directed Plus Per-Topology Stability Cache}

\textbf{Hypothesis.} Change two things at once. First, swap the
layered metaphor for a force-directed placement so that the
spatial layout reflects which nodes share connections rather than
imposing a left-to-right pipeline. Second, key the layout cache
by topology rather than by snapshot, so status and heartbeat
ticks do not trigger relayout at all.

\textbf{What shipped.} Commit \texttt{3c6d307} rewrote
\texttt{GraphTab.tsx} around three pieces, visible in the
current file:

\begin{itemize}[leftmargin=1.4em]
  \item \textbf{Force-directed layout via ELK's \texttt{force}
        algorithm} (the Eades model). The configuration sets
        \texttt{elk.algorithm} to \texttt{force},
        \texttt{elk.force.iterations} to 400, and uses a
        deterministic seed so the same topology produces the
        same layout (otherwise the cache would be useless).
  \item \textbf{Per-topology stability cache.} A
        \texttt{topologyKey} derived from the sorted list of node
        ids and edge ids drives the ELK effect. The
        \texttt{useEffect} that calls ELK depends only on
        \texttt{topologyKey}. Cached positions
        (\texttt{layout.positions}) are merged onto the current
        node set in a \texttt{useMemo}, so status and timestamp
        churn flow through into colours and labels without
        moving anything.
  \item \textbf{Per-side handles for smart edge routing.}
        Each node renders four target and four source handles
        (top, right, bottom, left). For each edge the
        lexicographically smaller endpoint picks the side
        facing its peer; the larger endpoint mirrors. Handle
        assignments are cached together with positions, so the
        same edge does not flip sides between ticks even if a
        cosmetic change runs through the
        \texttt{useMemo}. The \texttt{cancelled} flag on the
        layout promise drops stale ELK results from earlier
        ticks.
\end{itemize}

\textbf{Outcome.} The graph stays visually stable as agents and
topics come and go. A status change (online to degraded, for
example) recolours a node without moving it. The arrival of a
new agent shifts the layout once, when the new node's id enters
the \texttt{topologyKey}, and then stops. The operator can
follow what changed between ticks because what changed is the
only thing on the screen that moved.

% ==============================================================
\section{Method}
% ==============================================================

In rough order:

\begin{enumerate}[leftmargin=1.6em]
  \item Built the egui admission GUI alongside the router and
        ran it manually for a single-machine smoke test. The
        admission flow worked; the operator-surface gap was
        observable as soon as a second machine joined the test.
  \item Took Denis's launcher drop (\texttt{df4cbde}) and read
        it end-to-end before deciding to keep the Tauri shape
        and discard the egui shape. The decision was that the
        launcher's lifecycle (a desktop app the operator starts
        like any other) matched the platform's intended use; the
        router's lifecycle (a background process) did not.
  \item Pulled the router runtime up into a library
        (\texttt{1818090}). Made the egui dependency optional so
        the launcher can link the router without pulling a GUI
        toolkit it does not use. Re-derived the launcher's
        \texttt{InnerState} on top of the router's
        \texttt{AppState} so there is one source of truth for
        snapshots.
  \item Re-implemented the admission flow in React
        (\texttt{45cf40f}): chooser, requesting-join wait
        screen, pending-requests panel on the admin side. Kept
        the wire format on \texttt{dverse/session/} unchanged;
        the protocol from the security writeup did not move.
  \item Built the first graph prototype on dagre. Observed the
        layered metaphor flatten bidirectional pairs into
        back-edges in the running launcher.
  \item Swapped to ELK layered/orthogonal. Observed the edge
        routing improve and the node-jumping symptom persist.
        Wrote the finding down.
  \item Pivoted the layout strategy: force-directed placement,
        per-topology cache, per-side handles for smart edge
        routing. Verified that status ticks do not relayout the
        graph (positions stay byte-identical across snapshots
        that change only status fields) and that new-agent
        ticks relayout once and then settle.
\end{enumerate}

% ==============================================================
\section{Results}
% ==============================================================

Table~\ref{tab:ui-iter-result-shape} summarises the
operator-experience shape of each iteration. The entries are
qualitative, the way the security writeup's result table is
qualitative; the experiment is observational rather than
benchmarked.

\begin{table}[H]
  \centering
  \small
  \begin{tabular}{|p{4.6cm}|p{2.2cm}|p{2.6cm}|p{4cm}|}
    \hline
    \rowcolor{gray!15}
    \textbf{Property} & \textbf{V1 (egui)} & \textbf{V2 (Tauri)} & \textbf{V3 (graph)} \\
    \hline
    Operator windows open & 2 (router, plus whatever the rest of the UX lived in) & 1 & 1 \\
    \hline
    Joiner has a wait-state view & no & yes (\texttt{RequestingJoinScreen}) & yes \\
    \hline
    Admin admit/deny lives in & egui window & React panel inside main view & React panel inside main view \\
    \hline
    Graph view of agents and topics & no & no & yes \\
    \hline
    Graph stays stable across status ticks & N/A & N/A & yes (per-topology cache) \\
    \hline
  \end{tabular}
  \caption{Operator-experience shape of each iteration.}
  \label{tab:ui-iter-result-shape}
\end{table}

The end-to-end shape that V3 leaves behind is the launcher used
in the security writeup's two-machine demo: one window per
operator, login through chooser through main view, admission
panel on the admin's main view, graph tab one click away with
nodes that stop moving once the topology has settled. No
benchmarking was done; the claim that the graph is stable across
status ticks is from running the launcher and watching the
positions, the same way the security writeup's sub-second
agent-key relay claim is observational.

% ==============================================================
\section{What Remained Open}
% ==============================================================

V3 settles the core question (``can the launcher host the whole
operator surface without losing visual continuity as the network
changes?''), but leaves three items as follow-on work.

\textbf{Graph view of cross-session traffic.} The graph reflects
one session's agents and topics. Once V3 of the security model
allows multiple sessions on one fabric (see
\emph{Iterating on the Zenoh security model}, ``What Remained
Open''), the graph will either need a per-session view or a way
to overlay multiple sessions without making the picture
unreadable. Neither is built.

\textbf{Persisted layout across launcher restarts.} The
stability cache lives in component state. A launcher restart
loses it, and the first tick after restart relayouts the whole
graph. Persisting the cache to a small local store keyed by
topology would let restarts pick up where the last run left off.

\textbf{Operator-driven layout overrides.} A force-directed
layout sometimes places a node in a position the operator finds
inconvenient. There is no way to pin a node manually and have
ELK respect the pin on subsequent layouts. React Flow supports
drag-to-move; wiring the dragged position back into the cache
(and feeding it to ELK as a fixed coordinate) would close that
gap.

% ==============================================================
\section{Cross-References}
% ==============================================================

The admission protocol that \texttt{PendingRequestsPanel}
surfaces is defined in \emph{Iterating on the Zenoh security
model}: the JoinRequest, the admin's Allow or Deny decision, the
ACL coarsening and the Megolm key wrap that make admission
viable without a per-join data-plane outage. The router-as-
library refactor (\texttt{1818090}) that lets the launcher start
the router in-process is described in \emph{Framework
architecture and integration}, alongside the rest of the
\texttt{bot\_framework} composition story.

% ==============================================================
\section*{Glossary}
\addcontentsline{toc}{section}{Glossary}
% ==============================================================

\begin{description}[leftmargin=2.6cm, style=nextline, font=\sffamily\bfseries]
  \item[ACL] Access Control List. In Zenoh, a JSON block in
        the session configuration that defines which subjects
        are allowed or denied on which key expressions. Why
        ACLs cannot be the per-action admission knob is the
        subject of \emph{Iterating on the Zenoh security
        model}; what the launcher's pending-requests panel
        surfaces is the admission protocol that replaced the
        per-action ACL rewrite.
  \item[CN] Common Name. The subject field of a TLS leaf
        certificate. Used by the launcher's admission flow to
        identify a requester.
  \item[dagre] A JavaScript directed-graph layout library that
        implements the Sugiyama layered algorithm. Used in
        V3a, replaced in V3b.
  \item[eframe] The application framework around egui. The
        \texttt{RouterApp} struct in V1 implemented
        \texttt{eframe::App}.
  \item[egui] An immediate-mode Rust GUI library. The V1
        admission window was an egui window inside the
        \texttt{zenoh\_router} binary.
  \item[ELK] The Eclipse Layout Kernel. A Java layout engine
        with a JavaScript port (\texttt{elkjs}). Provides the
        layered, orthogonal, and force algorithms used in V3b
        and V3c.
  \item[force-directed layout] A graph-layout family in which
        edges act as springs pulling connected nodes together
        and nodes act as charged particles pushing each other
        apart. The placement reflects which nodes share
        connections rather than imposing a fixed reading
        direction. V3c uses ELK's \texttt{force} algorithm,
        which implements the Eades model.
  \item[JoinRequest] The signed message a candidate operator
        publishes to \path{dverse/session/requests/<cn>} to
        request admission. The launcher's
        \texttt{RequestingJoinScreen} is the joiner-side wait
        state for the matching admission decision.
  \item[mDNS] Multicast DNS. Used by the launcher's
        \texttt{discover\_routers} command to list sessions
        visible on the LAN for the chooser screen.
  \item[React] A JavaScript UI library. The launcher's
        frontend is a React application.
  \item[\texttt{rqt\_graph}] A ROS2 tool that draws nodes and
        topics as a graph. The V3 graph tab borrows its
        convention (ovals for nodes, rectangles for topics).
  \item[Tailwind] A utility-first CSS framework used by the
        launcher's frontend for component styling.
  \item[Tauri] A Rust framework for building cross-platform
        desktop applications with a web frontend and a Rust
        backend. V2 onward, the launcher is a Tauri 2 app.
  \item[topology] The set of agent nodes, topic nodes, and
        publishes / subscribes edges that the graph view
        renders. V3c's cache is keyed on this so the layout
        recomputes when (and only when) the topology changes.
  \item[Vite] The frontend build tool used by the launcher.
  \item[\texttt{@xyflow/react}] The React-Flow rendering
        library used by the graph tab to draw nodes, edges,
        and handles. The layout engine (dagre, then ELK) feeds
        \texttt{@xyflow/react} the coordinates.
  \item[Zenoh] A decentralised pub/sub messaging protocol from
        the Eclipse Foundation. The launcher embeds a Zenoh
        router as a library.
\end{description}

\end{document}
